\documentclass[UTF8,12pt,a4paper]{ctexart}

\usepackage{geometry}
\usepackage{booktabs}
\usepackage{longtable}
\usepackage{array}
\usepackage{tabularx}
\usepackage{xcolor}
\usepackage{hyperref}
\usepackage{enumitem}
\usepackage{listings}
\usepackage{amsmath}
\usepackage{graphicx}
\usepackage{caption}
\usepackage{float}

\geometry{left=2.4cm,right=2.4cm,top=2.4cm,bottom=2.4cm}
\hypersetup{
  colorlinks=true,
  linkcolor=blue,
  urlcolor=blue,
  citecolor=blue
}

\setlist[itemize]{leftmargin=2em}
\setlist[enumerate]{leftmargin=2em}
\renewcommand{\arraystretch}{1.25}

\lstdefinestyle{code}{
  basicstyle=\ttfamily\small,
  breaklines=true,
  frame=single,
  backgroundcolor=\color{gray!5},
  columns=fullflexible
}
\lstset{style=code}

\title{\textbf{TinyHear 双麦端侧降噪项目完整说明文档}\\
\large 面向助听器/耳戴设备上行链路的端侧 AI 降噪原型\\
\large Git 地址：\url{https://github.com/turambar928/TinyHear-Dual-Mic-Denoising}}
\author{}
\date{更新时间：2026 年 8 月 15 日}

\begin{document}

\maketitle
\tableofcontents
\newpage

\section{文档目的}

本文档用于向上级或项目评审人员完整说明 TinyHear 双麦端侧降噪项目的背景、目标、技术路线、实现细节、实验迭代、当前效果、已完成工程内容、现存问题和后续计划。读者不需要预先了解本项目代码，即可通过本文从零理解项目现状。

本项目不是单纯训练一个离线语音增强模型，而是围绕助听器/耳戴设备的端侧实时约束，构建一个从数据准备、模型训练、实时推理、网页试听、客观评估到部署验证的完整原型系统。

\section{项目背景}

\subsection{应用场景}

项目面向助听器、耳戴设备和上行通话场景。在这类设备中，用户讲话时设备需要把语音上传给手机或通信链路，因此上行音频质量直接影响通话可懂度和使用体验。

典型问题包括：

\begin{itemize}
  \item 环境噪声，例如街道、办公室、室内背景声；
  \item 风噪和气流噪声，例如走动、户外风、佩戴设备周围气流；
  \item 麦克风自噪声和摩擦噪声；
  \item 近距双麦条件下空间差异弱，传统简单双麦处理效果有限；
  \item 设备算力、内存、功耗和延迟都受限。
\end{itemize}

\subsection{为什么使用双麦}

相较单麦，双麦系统可以利用两个麦克风之间的空间差异，例如相位差、相干性和局部空间统计信息。对于上行降噪，双麦信息有助于区分目标语音和空间上不同的噪声源。

但助听器/耳戴设备上的麦克风距离通常很近，两个麦克风之间的幅度差（ILD, inter-channel level difference）不明显。因此，本项目后续重点使用：

\begin{itemize}
  \item IPD（inter-channel phase difference，相位差）；
  \item coherence（双麦相干性）；
  \item coherence-weighted MWF 风格空间前端；
  \item 多帧/复数域滤波，而不是只做单帧幅度 mask。
\end{itemize}

\section{项目初始目标与约束}

项目初始要求可以概括为：

\begin{quote}
做一个双麦克风上行降噪模型，模型大小控制在 100K--150K 参数左右，模型架构可自由设计，数据集可以先用公开数据集，最好能用整型运算完成。
\end{quote}

当前项目与这些要求的对应关系如表 \ref{tab:requirements} 所示。

\begin{table}[H]
\centering
\caption{项目要求与当前状态}
\label{tab:requirements}
\begin{tabularx}{\textwidth}{p{3.8cm}p{3.0cm}X}
\toprule
\textbf{要求} & \textbf{当前状态} & \textbf{说明} \\
\midrule
双麦上行降噪 & 已满足 & 当前输入为双通道麦克风音频，输出单通道增强语音，并显式使用双麦空间特征。 \\
模型大小 100K--150K & 已满足 & 当前推荐模型 TinyDeepFilterTCN 参数量为 137,984。 \\
架构可自由设计 & 已完成多轮探索 & 已尝试 band-mask TCN、空间特征 TCN、DeepFilter、GRU、complex mask、MWF student、MWF teacher 蒸馏等路线。 \\
公开数据集 & 已满足 & 当前主要使用 CMU ARCTIC、DEMAND 和 Zenodo Wind Noise Dataset。 \\
整型运算 & 部分完成 & 早期 band-mask TCN 已做 INT8/C reference 验证；当前最佳 DeepFilter 版本仍需完成 INT8 导出和 C reference 对齐。 \\
\bottomrule
\end{tabularx}
\end{table}

\section{当前推荐版本概览}

当前推荐版本为：

\begin{lstlisting}
teacher_deepfilter_conservative_on_wind_eval
\end{lstlisting}

对应 checkpoint：

\begin{lstlisting}
runs/arctic_wind_teacher_deepfilter_mwf_conservative/best.pt
\end{lstlisting}

该版本的核心方案是：

\begin{lstlisting}
coherence/MWF 双麦空间前端
  + Tiny DeepFilter TCN
  + conservative MWF teacher 蒸馏
  + 风噪数据 fine-tune
\end{lstlisting}

需要特别说明的是，版本名中的 \texttt{teacher} 并不表示实际部署时运行 teacher 模型。Teacher 是训练阶段使用的 oracle local MWF 目标，用于指导小模型学习更接近强双麦滤波器的行为。实际推理和网页 demo 中的 Realtime 音频，运行的是 137,984 参数的 TinyDeepFilterTCN 小模型。

\section{系统整体流程}

\subsection{推理流程}

当前主线推理流程如下：

\begin{lstlisting}
双麦 waveform
  -> STFT
  -> IPD / coherence 等空间特征
  -> coherence-weighted MWF spatial frontend
  -> Tiny DeepFilter TCN
  -> band gain + low-frequency deep filtering
  -> iSTFT / overlap-add
  -> enhanced speech
\end{lstlisting}

推理时没有 clean 语音，也没有 oracle teacher。模型只接收双麦 mixture，并输出增强后的单通道语音。

\subsection{展示流程}

网页 demo 路径：

\begin{lstlisting}
runs/audio_demo/index.html
\end{lstlisting}

常用访问地址：

\begin{lstlisting}
http://127.0.0.1:38180/runs/audio_demo/index.html
\end{lstlisting}

网页中每个样本通常包含：

\begin{itemize}
  \item \textbf{Noisy}：输入噪声音频；
  \item \textbf{Realtime / Enhanced}：当前模型实时链路输出，代表实际部署方向的效果；
  \item \textbf{Clean}：干净语音参考，只用于对比，不参与实际推理。
\end{itemize}

\section{模型结构}

\subsection{当前推荐模型}

当前推荐模型为 \texttt{TinyDeepFilterTCN}，主要配置如表 \ref{tab:model_config} 所示。

\begin{table}[H]
\centering
\caption{当前推荐模型配置}
\label{tab:model_config}
\begin{tabular}{ll}
\toprule
\textbf{项目} & \textbf{数值} \\
\midrule
模型类型 & TinyDeepFilterTCN \\
参数量 & 137,984 \\
采样率 & 16 kHz \\
STFT & 256 点 FFT \\
hop length & 64 samples，即 4 ms \\
channels & 96 \\
blocks & 8 \\
kernel size & 5 \\
df bins & 64 \\
df order & 3 \\
spatial frontend & coherence\_mwf \\
\bottomrule
\end{tabular}
\end{table}

\subsection{模型输入}

模型输入不是单麦幅度谱，而是结合了：

\begin{itemize}
  \item 参考通道时频特征；
  \item 双麦相位差；
  \item 双麦相干性；
  \item coherence/MWF 前端输出的参考信息。
\end{itemize}

这样做的目的，是避免模型退化成普通单麦降噪，让它能利用近距双麦中的空间信息。

\subsection{模型输出}

模型输出两类结果：

\begin{enumerate}
  \item \textbf{band gain}：用于整体谱幅度抑制，稳定且便于端侧实现；
  \item \textbf{deep-filter complex coefficients}：主要作用于低频区域，用多帧复数滤波缓解只做幅度 mask 带来的相位问题和低频风噪残留。
\end{enumerate}

\section{数据集与数据构造}

\subsection{使用的数据集}

项目先后使用的数据集如表 \ref{tab:datasets} 所示。

\begin{table}[H]
\centering
\caption{数据集演进}
\label{tab:datasets}
\begin{tabularx}{\textwidth}{p{3.2cm}p{3.2cm}p{3.5cm}X}
\toprule
\textbf{阶段} & \textbf{Clean speech} & \textbf{Noise} & \textbf{用途} \\
\midrule
synthetic baseline & 合成语音 & 合成噪声 & 快速打通训练和推理流程 \\
YESNO + DEMAND & OpenSLR YESNO & DEMAND & 小规模公开数据 sanity check \\
ARCTIC + DEMAND & CMU ARCTIC & DEMAND & 主基线 \\
ARCTIC + DEMAND + wind & CMU ARCTIC & DEMAND + Zenodo Wind Noise & 当前风噪优化主线 \\
\bottomrule
\end{tabularx}
\end{table}

当前主线训练目录：

\begin{lstlisting}
data/arctic_wind_demand_flat/
\end{lstlisting}

固定评估集：

\begin{lstlisting}
data/arctic_wind_eval/val
data/arctic_demand_eval/val
\end{lstlisting}

原始下载数据和整理后的大数据目录不提交到 GitHub；仓库中主要保留代码、文档、checkpoint、metrics、少量 demo 音频和分析报告。

\subsection{双麦数据构造}

公开 clean speech 与 noise 数据本身不一定天然符合助听器双麦形式，因此项目中构造了双麦 mixture。大致逻辑包括：

\begin{itemize}
  \item clean speech 归一化到固定 RMS；
  \item 按随机 SNR 混入 noise；
  \item 根据随机入射方向和麦克风距离模拟双麦时间差；
  \item 采样不同近距双麦距离，当前范围约为 0.015--0.024 m；
  \item 加入自噪声和风噪增强；
  \item 输出双通道 mixture 和对应 clean pair。
\end{itemize}

\section{MWF Teacher 蒸馏}

\subsection{为什么引入 Teacher}

早期版本虽然能降噪，但网页试听中反复暴露出两个问题：

\begin{itemize}
  \item 底噪和低频 ``呼呼声'' 仍明显；
  \item 只靠 postfilter、dehiss、gate 或更换小模型，收益有限。
\end{itemize}

后续用 oracle local MWF 做上限验证时，发现强双麦前端在风噪评估集上明显更强：

\begin{lstlisting}
Oracle local MWF on Zenodo wind eval:
Noisy SI-SDR:     5.39 dB
Enhanced SI-SDR: 14.21 dB
Improvement:     +8.82 dB
\end{lstlisting}

这说明强双麦前端/MWF 是有效方向。但 oracle MWF 依赖 clean/noise 信息，无法在真实部署时直接使用。因此，本项目把 oracle local MWF 作为训练阶段的 teacher，让小模型学习其滤波行为。

\subsection{Teacher 与部署模型的区别}

训练阶段可以拿到：

\begin{itemize}
  \item noisy 双麦；
  \item clean 双麦。
\end{itemize}

因此可以构造 oracle local MWF teacher。实际部署时只有 noisy 双麦，没有 clean，因此 teacher 不会参与推理。

可以概括为：

\begin{lstlisting}
训练阶段：noisy + clean -> oracle MWF teacher -> 指导 TinyDeepFilterTCN
部署阶段：noisy 双麦 -> TinyDeepFilterTCN -> enhanced speech
\end{lstlisting}

\subsection{保守 Teacher 目标}

项目中也尝试过 aggressive teacher，即让小模型大比例模仿 oracle MWF。但实验发现 teacher 太强会导致小模型学习不稳定，出现过度滤波、人声变小和通用噪声场景退化。

因此当前推荐版本使用 conservative teacher：

\begin{equation}
\text{target} = 0.25 \cdot \text{oracle\_mwf\_teacher} + 0.75 \cdot \text{clean\_mic0}
\end{equation}

关键参数：

\begin{lstlisting}
teacher_blend: 0.25
teacher_mask_weight: 0.35
oracle_window: 9
diagonal_loading: 0.01
\end{lstlisting}

\section{Loss 设计}

当前训练没有只使用 MSE 或单一 SI-SDR，而是组合多个目标，以同时约束语音质量、谱幅度、残留噪声和滤波稳定性。

\begin{table}[H]
\centering
\caption{当前训练 loss 组成}
\label{tab:loss}
\begin{tabularx}{\textwidth}{p{4.2cm}X}
\toprule
\textbf{Loss} & \textbf{作用} \\
\midrule
masked MSE & 约束预测 band gain 与目标 mask 的差异 \\
waveform L1 loss & 约束增强波形与目标波形接近 \\
STFT log magnitude loss & 约束听感相关的谱幅度结构 \\
SI-SDR loss & 提升整体语音重建质量 \\
coef energy loss & 限制 deep-filter 复数系数过大，减少不稳定伪影 \\
residual noise loss & 对 clean 中能量较低的频点施加残留噪声惩罚 \\
silence floor loss & 压低静音和弱语音区域的输出底噪 \\
\bottomrule
\end{tabularx}
\end{table}

当前保守 teacher 版本关键权重：

\begin{lstlisting}
waveform_loss_weight: 0.70
stft_mag_loss_weight: 0.22
si_sdr_loss_weight: 0.04
coef_reg_weight: 0.002
residual_noise_loss_weight: 0.34
silence_floor_weight: 0.28
\end{lstlisting}

\section{版本迭代过程}

\subsection{初始 Tiny band-mask TCN}

最早版本使用轻量 causal TCN 输出 32-band mask。优点是模型小、易于 INT8、易于 C reference，也能实现基本降噪。

问题是：

\begin{itemize}
  \item 人声容易变小；
  \item 底噪明显；
  \item 只做幅度 mask，无法处理 noisy phase 和细粒度残留噪声。
\end{itemize}

\subsection{空间特征增强}

后续加入 IPD、coherence 和双麦能量相关特征。原因是近距双麦下 ILD 差异较小，相位差和相干性更有价值。该阶段让模型真正利用双麦输入，而不是退化为单麦降噪。

\subsection{Gate 与 Bypass}

尝试使用 learned gate 或 high-SNR bypass，减少干净或高 SNR 场景下的过处理。这对部分样本有帮助，但不能根治底噪和风噪呼呼声。

\subsection{Tiny DeepFilter}

将后端从简单 band mask 升级为 Tiny DeepFilter。DeepFilter 不只输出频带增益，还输出低频复数滤波系数，可以做多帧复数域滤波。

该路线让客观指标和主观听感明显优于早期 band-mask，但仍存在低频风噪残留。

\subsection{coherence/MWF 空间前端}

前端从简单 delay-sum 升级为 coherence-weighted MWF 风格空间前端，用双麦相干性生成更好的参考通道，再交给 DeepFilter 后端细化。

这一方向成为后续主线。

\subsection{后处理、GRU、Complex Mask 与 Student 分支}

项目中还尝试过：

\begin{itemize}
  \item stable postfilter；
  \item dehiss postfilter；
  \item airflow residual filter；
  \item RNNoise-style TinyGRU；
  \item full-bin complex mask TCN；
  \item MWF covariance student。
\end{itemize}

这些方法对局部指标或部分样本有改善，但主观听感不稳定，没有明显超过 DeepFilter + 强双麦前端主线。因此当前主线回到 Tiny DeepFilter + coherence/MWF frontend + conservative MWF teacher。

\section{当前实验结果}

\subsection{当前推荐 checkpoint}

\begin{lstlisting}
runs/arctic_wind_teacher_deepfilter_mwf_conservative/best.pt
\end{lstlisting}

\subsection{客观指标}

当前推荐版本在两个固定评估集上的结果如表 \ref{tab:metrics} 所示。

\begin{table}[H]
\centering
\caption{当前推荐版本客观指标}
\label{tab:metrics}
\begin{tabular}{lrrrr}
\toprule
\textbf{评估集} & \textbf{items} & \textbf{Noisy SI-SDR} & \textbf{Enhanced SI-SDR} & \textbf{Improvement} \\
\midrule
Zenodo wind eval & 160 & 5.39 dB & 10.19 dB & +4.80 dB \\
Original DEMAND eval & 160 & 4.33 dB & 10.96 dB & +6.63 dB \\
\bottomrule
\end{tabular}
\end{table}

输出/输入 RMS ratio：

\begin{itemize}
  \item wind eval：0.764；
  \item DEMAND eval：0.690。
\end{itemize}

这些结果说明当前模型具备明确降噪能力。但从网页 demo 主观听感看，相比 clean 仍存在低频风噪/气流声残留。

\subsection{关键版本对比}

\begin{table}[H]
\centering
\caption{关键版本对比}
\label{tab:version_compare}
\begin{tabularx}{\textwidth}{p{4.0cm}p{2.2cm}p{2.2cm}X}
\toprule
\textbf{版本} & \textbf{Wind eval} & \textbf{DEMAND eval} & \textbf{结论} \\
\midrule
wind fine-tune & +4.35 dB & +7.06 dB & 通用噪声更稳，但风噪仍有明显残留。 \\
aggressive teacher & +3.74 dB & +5.21 dB & teacher 太强，小模型学不稳，出现退化。 \\
conservative teacher & +4.80 dB & +6.63 dB & 当前推荐，风噪指标和听感更均衡。 \\
\bottomrule
\end{tabularx}
\end{table}

当前版本不是所有指标都最高，而是在风噪抑制、人声保持和模型稳定性之间更均衡。

\section{失败样本分析}

为避免只依赖主观听感，项目新增失败模式分析脚本：

\begin{lstlisting}
scripts/analyze_failure_modes.py
\end{lstlisting}

wind eval 的失败模式统计如表 \ref{tab:failure} 所示。

\begin{table}[H]
\centering
\caption{Wind eval 失败模式统计}
\label{tab:failure}
\begin{tabular}{lr}
\toprule
\textbf{Failure mode} & \textbf{数量} \\
\midrule
ok & 117 \\
low\_freq\_wind & 35 \\
regression & 7 \\
residual\_noise & 1 \\
\bottomrule
\end{tabular}
\end{table}

关键均值：

\begin{itemize}
  \item Mean SI-SDR improvement：+4.80 dB；
  \item Mean speech preservation：-0.96 dB；
  \item Mean quiet noise reduction：15.73 dB；
  \item Mean quiet enhanced RMS：-44.32 dBFS。
\end{itemize}

这个分析说明，当前主要问题不是人声整体被严重压小，而是部分样本存在低频风噪/气流残留。因此后续不应该盲目提高全频抑制强度，而应该围绕 \texttt{low\_freq\_wind} 样本做定向优化。

\section{工程实现与代码结构}

项目核心目录如下：

\begin{lstlisting}
src/ha_denoise/                 # 核心模型、特征、数据和空间处理
scripts/                        # 训练、评估、增强、demo 和分析脚本
c_reference/                    # 早期 C reference 和 benchmark
runs/                           # checkpoint、metrics、demo 音频和报告
docs/                           # 项目文档、汇报材料和说明
\end{lstlisting}

关键文件：

\begin{table}[H]
\centering
\caption{关键代码文件}
\label{tab:files}
\begin{tabularx}{\textwidth}{p{5.2cm}X}
\toprule
\textbf{文件} & \textbf{作用} \\
\midrule
\texttt{src/ha\_denoise/model.py} & TinyCausalTCN、TinyDeepFilterTCN 等模型结构 \\
\texttt{src/ha\_denoise/features.py} & STFT、频带特征和时频处理 \\
\texttt{src/ha\_denoise/spatial.py} & 双麦空间特征、coherence/MWF 前端相关逻辑 \\
\texttt{src/ha\_denoise/dataset.py} & 动态双麦混音数据集构造 \\
\texttt{scripts/train\_deepfilter.py} & 当前 DeepFilter 主线训练脚本，包含 MWF teacher 蒸馏 \\
\texttt{scripts/evaluate\_deepfilter.py} & DeepFilter 评估脚本 \\
\texttt{scripts/enhance\_realtime.py} & 实时增强链路脚本 \\
\texttt{scripts/build\_audio\_demo.py} & 网页试听 demo 生成脚本 \\
\texttt{scripts/analyze\_failure\_modes.py} & 失败模式分析脚本 \\
\bottomrule
\end{tabularx}
\end{table}

\section{如何测试和试听项目}

本节面向第一次接触项目的使用者，说明从 GitHub 拉取项目后，如何在不重新训练的情况下试听不同模型版本的效果，以及如何用命令行对单条音频运行当前推荐模型。

\subsection{拉取项目}

在一台已经安装 Git 和 Python 3 的机器上执行：

\begin{lstlisting}
git clone https://github.com/turambar928/TinyHear-Dual-Mic-Denoising.git
cd TinyHear-Dual-Mic-Denoising
\end{lstlisting}

如果在服务器上已经有项目目录，则直接进入项目：

\begin{lstlisting}
cd /home/taozifu2025/internship
\end{lstlisting}

当前仓库已经提交了网页 demo、少量试听样本、关键 checkpoint 和 metrics。因此，如果只是想听不同模型版本的效果，不需要重新下载完整训练数据，也不需要重新训练。

\subsection{安装运行依赖}

建议先创建虚拟环境：

\begin{lstlisting}
python3 -m venv .venv
source .venv/bin/activate
pip install -r requirements.txt
\end{lstlisting}

项目基础依赖包括：

\begin{itemize}
  \item numpy；
  \item scipy；
  \item torch；
  \item tqdm。
\end{itemize}

如果只是打开网页听已经生成好的 wav 文件，严格来说不需要安装 PyTorch；但如果要运行模型推理、重新评估或生成新的 enhanced wav，则需要安装 \texttt{requirements.txt} 中的依赖。

\subsection{最快试听方式：打开网页 Demo}

进入项目根目录后启动一个静态 HTTP 服务：

\begin{lstlisting}
python3 -m http.server 38180 --bind 0.0.0.0
\end{lstlisting}

然后在浏览器打开：

\begin{lstlisting}
http://127.0.0.1:38180/runs/audio_demo/index.html
\end{lstlisting}

如果是在 VSCode 远程服务器中运行，需要在 VSCode 的 Ports/端口面板中转发 38180 端口。若手动使用 SSH 跳板机转发，可以在本地电脑执行：

\begin{lstlisting}
ssh -N -L 15180:192.168.14.16:38180 taozifu2025@124.16.139.208 -p 17022
\end{lstlisting}

然后在本地浏览器打开：

\begin{lstlisting}
http://127.0.0.1:15180/runs/audio_demo/index.html
\end{lstlisting}

如果本地 15180 被占用，可以换成其他本地端口，例如 15181：

\begin{lstlisting}
ssh -N -L 15181:192.168.14.16:38180 taozifu2025@124.16.139.208 -p 17022
\end{lstlisting}

对应浏览器地址改为：

\begin{lstlisting}
http://127.0.0.1:15181/runs/audio_demo/index.html
\end{lstlisting}

\subsection{网页里应该听哪些版本}

网页 demo 中会列出多个历史版本，每个版本下面包含若干样本。每个样本通常有四个播放器：

\begin{itemize}
  \item \textbf{Noisy}：输入噪声音频；
  \item \textbf{Offline}：离线增强结果；
  \item \textbf{Realtime}：按实时链路生成的增强结果；
  \item \textbf{Clean}：干净语音参考。
\end{itemize}

推荐优先听：

\begin{lstlisting}
teacher_deepfilter_conservative_on_wind_eval
\end{lstlisting}

这个版本对应当前推荐部署候选模型：

\begin{lstlisting}
runs/arctic_wind_teacher_deepfilter_mwf_conservative/best.pt
\end{lstlisting}

试听顺序建议如下：

\begin{enumerate}
  \item 先播放 \textbf{Noisy}，确认输入中的风噪或背景噪；
  \item 再播放 \textbf{Realtime}，这是当前小模型实际推理链路的效果；
  \item 最后播放 \textbf{Clean}，用于理解当前模型和理想目标之间的差距。
\end{enumerate}

为了理解版本迭代，可以重点对比表 \ref{tab:demo_versions} 中的版本。

\begin{table}[H]
\centering
\caption{网页 demo 中建议对比的版本}
\label{tab:demo_versions}
\begin{tabularx}{\textwidth}{p{5.2cm}X}
\toprule
\textbf{网页版本名} & \textbf{含义和用途} \\
\midrule
\texttt{teacher\_deepfilter\_conservative\_on\_wind\_eval} & 当前推荐版本，风噪场景下综合最好；实际听 Realtime。 \\
\texttt{teacher\_deepfilter\_conservative\_on\_demand\_eval} & 同一模型在 DEMAND 通用噪声评估集上的效果。 \\
\texttt{wind\_finetune\_on\_wind\_eval} & 未加入 conservative teacher 前的风噪 fine-tune 版本，用于比较 teacher 蒸馏收益。 \\
\texttt{teacher\_deepfilter\_aggressive\_on\_wind\_eval} & 激进 teacher 版本，说明 teacher 目标太强时可能导致退化。 \\
\texttt{tiny\_deepfilter\_coherence\_mwf} & 早期 DeepFilter + coherence/MWF 主线，可用于对比风噪数据和 teacher 之前的效果。 \\
\texttt{baseline} & 最早 band-mask TCN 基线，用于展示从初始模型到当前模型的变化。 \\
\bottomrule
\end{tabularx}
\end{table}

汇报或验收时，建议不要只播放一个增强结果，而是按照：

\begin{lstlisting}
Noisy -> Realtime -> Clean
\end{lstlisting}

的顺序播放，这样听众能更清楚地区分“已有降噪能力”和“距离 clean 仍有差距”。

\subsection{如何确认网页服务是否正常}

如果浏览器打不开网页，先在服务器上检查 HTTP 服务是否正常：

\begin{lstlisting}
curl -I --max-time 3 http://127.0.0.1:38180/runs/audio_demo/index.html
\end{lstlisting}

如果返回类似：

\begin{lstlisting}
HTTP/1.0 200 OK
\end{lstlisting}

说明服务器上的网页服务和文件路径没有问题。此时打不开通常是本地电脑到服务器的端口转发问题，需要检查 VSCode Ports 或 SSH \texttt{-L} 转发命令。

也可以检查 38180 端口是否有进程监听：

\begin{lstlisting}
ss -ltnp | grep 38180
\end{lstlisting}

\subsection{命令行生成单条增强音频}

如果想不用网页，而是对某条 wav 直接运行当前推荐模型，可以执行：

\begin{lstlisting}
PYTHONPATH=src python3 scripts/enhance_realtime.py \
  --checkpoint runs/arctic_wind_teacher_deepfilter_mwf_conservative/best.pt \
  --input data/arctic_wind_eval/val/mix_0000.wav \
  --output runs/arctic_wind_teacher_deepfilter_mwf_conservative/test_realtime.wav \
  --loudness-match
\end{lstlisting}

输出文件为：

\begin{lstlisting}
runs/arctic_wind_teacher_deepfilter_mwf_conservative/test_realtime.wav
\end{lstlisting}

注意，不能在 shell 里直接输入 wav 文件路径来播放音频。例如：

\begin{lstlisting}
runs/arctic_wind_teacher_deepfilter_mwf_conservative/test_realtime.wav
\end{lstlisting}

这样会被 shell 当作可执行文件，因此可能出现“权限不够”。正确做法是用网页播放器、系统音频播放器，或使用命令行播放器，例如 \texttt{aplay}、\texttt{ffplay} 等。

\subsection{重新生成网页 Demo}

如果后续又生成了新的 listening eval 目录，可以用脚本重建网页：

\begin{lstlisting}
PYTHONPATH=src python3 scripts/build_audio_demo.py \
  --out runs/audio_demo/index.html
\end{lstlisting}

该脚本默认会收集项目中配置好的多个历史版本，并生成统一网页。新增版本也可以通过脚本的 \texttt{--set} 参数加入，格式为：

\begin{lstlisting}
PYTHONPATH=src python3 scripts/build_audio_demo.py \
  --set my_new_model=runs/my_new_model/listening_eval \
  --out runs/audio_demo/index.html
\end{lstlisting}

\subsection{重新评估一个 checkpoint}

如果要重新跑当前推荐模型在 wind eval 上的客观指标和试听样本，可以使用：

\begin{lstlisting}
PYTHONPATH=src python3 scripts/evaluate_deepfilter.py \
  --checkpoint runs/arctic_wind_teacher_deepfilter_mwf_conservative/best.pt \
  --data data/arctic_wind_eval \
  --split val \
  --save-listening runs/arctic_wind_teacher_deepfilter_mwf_conservative/listening_eval_wind_zenodo \
  --listening-samples 5 \
  --dehiss-postfilter \
  --dehiss-strength 1.05 \
  --dehiss-low-floor 0.82 \
  --dehiss-high-floor 0.42 \
  --dehiss-high-start-hz 2600 \
  --dehiss-full-strength-hz 5200
\end{lstlisting}

重新评估需要本地存在对应数据目录。如果只是查看当前项目进展，优先使用已经生成好的网页 demo。

\section{训练与评估命令}

当前推荐版本训练命令核心如下：

\begin{lstlisting}
PYTHONPATH=src python3 scripts/train_deepfilter.py \
  --data data/arctic_wind_demand_flat \
  --out runs/arctic_wind_teacher_deepfilter_mwf_conservative \
  --resume runs/arctic_wind_demand_tiny_deepfilter_coherence_mwf/best.pt \
  --reset-best-on-resume \
  --epochs 4 \
  --batch-size 8 \
  --seconds 1.0 \
  --on-the-fly \
  --virtual-multiplier 2 \
  --snr-min-db -12 \
  --snr-max-db 8 \
  --noise-mix-prob 0.65 \
  --mic-distance-min-m 0.015 \
  --mic-distance-max-m 0.024 \
  --self-noise-prob 0.45 \
  --self-noise-db -42 \
  --wind-noise-prob 0.10 \
  --wind-noise-db -32 \
  --spatial-frontend coherence_mwf \
  --teacher-mwf \
  --teacher-mask-weight 0.35 \
  --teacher-blend 0.25 \
  --oracle-window 9 \
  --diagonal-loading 0.01 \
  --channels 96 \
  --blocks 8 \
  --kernel-size 5 \
  --df-bins 64 \
  --df-order 3 \
  --waveform-loss-weight 0.70 \
  --stft-mag-loss-weight 0.22 \
  --si-sdr-loss-weight 0.04 \
  --coef-reg-weight 0.002 \
  --residual-noise-loss-weight 0.34 \
  --silence-floor-weight 0.28 \
  --min-gain 0.005 \
  --lr 1e-4 \
  --device cuda
\end{lstlisting}

网页 demo 启动方式：

\begin{lstlisting}
cd /home/taozifu2025/internship
python3 -m http.server 38180 --bind 0.0.0.0
\end{lstlisting}

浏览器访问：

\begin{lstlisting}
http://127.0.0.1:38180/runs/audio_demo/index.html
\end{lstlisting}

如果使用远程服务器，需要确认 VSCode 端口转发或 SSH port forwarding 已经把本地端口转发到服务器的 38180。

\section{当前成果}

截至目前，项目已经完成：

\begin{itemize}
  \item 双麦端侧降噪项目定位和初始需求拆解；
  \item 公开数据集准备，包括 CMU ARCTIC、DEMAND 和 Zenodo Wind Noise；
  \item 动态双麦混音训练集构造；
  \item Tiny band-mask TCN baseline；
  \item Tiny DeepFilter TCN 主线模型；
  \item IPD/coherence 双麦空间特征；
  \item coherence/MWF 空间前端；
  \item oracle local MWF teacher 蒸馏训练；
  \item 风噪数据 fine-tune；
  \item 多版本训练、评估和主观试听对比；
  \item 网页 demo，可直接听 Noisy、Realtime、Clean；
  \item failure analysis，用于定位低频风噪、回退、人声变小等问题；
  \item 早期 INT8 导出和 C reference 验证链路；
  \item README、项目报告、PPT 和汇报文档。
\end{itemize}

\section{当前不足}

当前项目仍存在以下问题：

\begin{enumerate}
  \item \textbf{和 clean 仍有听感差距}：当前版本具备降噪能力，但部分样本仍有低频风噪和气流呼呼声。
  \item \textbf{真实数据不足}：当前主要是公开 clean speech + noise 合成双麦数据，与真实助听器佩戴位置、麦克风自噪声、风噪和摩擦噪声存在 domain gap。
  \item \textbf{风噪和通用噪声存在取舍}：conservative teacher 版本提升 wind eval，但 DEMAND eval 相比普通 wind fine-tune 略有下降。
  \item \textbf{主观听感与 SI-SDR 不完全一致}：部分版本指标不错，但试听效果并不一定更好。
  \item \textbf{当前最佳模型尚未完成完整整型部署闭环}：早期 band-mask TCN 已验证 INT8/C reference，但当前 TinyDeepFilterTCN 仍需完成导出、定点参考和 C 端对齐。
\end{enumerate}

\section{下一步计划}

建议后续按以下优先级推进。

\subsection{低频风噪定向优化}

使用 failure analysis 中标记为 \texttt{low\_freq\_wind} 的样本作为固定失败集，重点优化低频风噪残留。

具体方向：

\begin{itemize}
  \item 增加低频 residual wind loss；
  \item 对静音/弱语音区域低频残留加约束；
  \item 避免全频强压噪导致人声损伤。
\end{itemize}

\subsection{语音保持约束}

继续压低风噪时，需要防止人声再次变小。

方向：

\begin{itemize}
  \item speech-active RMS preservation；
  \item speech-band loss；
  \item 高 SNR 场景减少过处理。
\end{itemize}

\subsection{真实双麦数据补充}

如果要进一步接近真实助听器效果，需要采集或引入更真实的双麦数据：

\begin{itemize}
  \item 佩戴位置双麦语音；
  \item 真实风噪；
  \item 走动气流；
  \item 摩擦噪声；
  \item 不同房间混响；
  \item 设备麦克风自噪声。
\end{itemize}

即使数据量较小，也可以作为 fine-tune 数据或固定 failure set。

\subsection{端侧部署闭环}

当前推荐 TinyDeepFilterTCN 需要继续完成：

\begin{itemize}
  \item INT8 权重导出；
  \item fixed-scale Python reference；
  \item C reference 对齐；
  \item 每帧耗时、状态内存和模型大小 benchmark；
  \item 后续适配 CMSIS-DSP/CMSIS-NN。
\end{itemize}

\subsection{展示和评估完善}

\begin{itemize}
  \item 网页 demo 保持推荐版本在最前；
  \item 增加固定风噪失败样本；
  \item 汇报时同时展示 SI-SDR、failure analysis 和主观试听；
  \item 建立固定听感评价流程。
\end{itemize}

\section{阶段性结论}

本项目已经完成从初始轻量降噪模型到当前双麦 MWF teacher + Tiny DeepFilter 主线的完整迭代。当前推荐版本满足双麦输入、100K--150K 参数、小模型实时推理、公开数据训练和网页 demo 展示等核心要求。

最重要的阶段性结论是：

\begin{quote}
项目现在的主要瓶颈不是 ``有没有降噪能力''，而是低频风噪残留、真实双麦数据匹配度，以及当前最佳模型的整型部署闭环。
\end{quote}

因此，后续不应无目的地频繁切换模型结构，而应围绕当前最有效主线继续收敛：

\begin{lstlisting}
coherence/MWF spatial frontend
  + Tiny DeepFilter TCN
  + conservative MWF teacher distillation
\end{lstlisting}

下一阶段重点是低频风噪定向优化、真实双麦数据补充和当前 DeepFilter 主线的 INT8/C reference 部署验证。

\appendix


\section{附录：关键路径}

\begin{lstlisting}
README.md
docs/current_progress_summary.md
docs/conservative_mwf_teacher_deepfilter.md
docs/optimization_history_report.md
docs/teacher_meeting_report.md
docs/presentation/tinyhear_project_intro.pptx
runs/audio_demo/index.html
runs/arctic_wind_teacher_deepfilter_mwf_conservative/best.pt
runs/arctic_wind_teacher_deepfilter_mwf_conservative/eval_wind_zenodo/metrics.json
runs/arctic_wind_teacher_deepfilter_mwf_conservative/analysis_wind_zenodo/failure_analysis.md
\end{lstlisting}

\end{document}
